% Part: first-order-logic
% Chapter: model-theory
% Section: partial-iso

\documentclass[../../../include/open-logic-section]{subfiles}

\begin{document}

\olfileid{mod}{bas}{dlo}
\section{Dichte lineaire ordeningen: altijd ruimte ertussen}

\begin{defn}
  Een \emph{dichte lineaire ordening zonder eindpunten} is een !!a{structure}
  $\Struct{M}$ met één manier om twee elementen te vergelijken. De
  !!{language} heeft daarvoor één tweeplaatsig !!{predicate}~$<$. De
  structuur moet de volgende zinnen waar maken:
  \begin{enumerate}
  \item $\lforall[x][\lnot x < x]$: niets staat vóór zichzelf;
  \item $\lforall[x][\lforall[y][\lforall[z][(x < y \lif (y < z \lif x
    <z ))]]]$: de volgorde loopt door;
  \item $\lforall[x][\lforall[y][(x< y \lor \eq[x][y] \lor y < x)]]$:
    twee elementen zijn gelijk of het ene staat vóór het andere;
  \item $\lforall[x][\lexists[y][x < y]]$: na elk element komt nog een element;
  \item $\lforall[x][\lexists[y][y < x]]$: vóór elk element komt nog een element;
  \item $\lforall[x][\lforall[y][(x < y \lif \lexists[z][(x < z \land
        z < y))]]]$: tussen twee elementen die na elkaar staan, zit nog een element.
 \end{enumerate}
\end{defn}

\begin{thm}\ollabel{thm:cantorQ}
  Neem twee !!{enumerable} dichte lineaire ordeningen zonder
  eindpunten. Ze zijn isomorf: de elementen kunnen zo aan elkaar worden
  gekoppeld dat de hele ordening hetzelfde blijft.
\end{thm}

\begin{proof}
  Neem $\Struct{M_1}$ en $\Struct{M_2}$ als twee !!{enumerable} dichte lineaire
  ordeningen zonder eindpunten. Schrijf ${<_1} = \Assign{<}{M_1}$ en ${<_2} =
  \Assign{<}{M_2}$. Laat $\PIso{I}$ de verzameling zijn van alle partiële
  isomorfismen tussen de twee ordeningen. De verzameling $\PIso{I}$ is niet leeg, want
  in elk geval $\emptyset \in \PIso{I}$. We laten zien dat je elk van die
  gedeeltelijk ingevulde koppelingen aan beide kanten kunt uitbreiden: $\PIso{I}$
  voldoet aan de heen-en-weereigenschap. Dan $\Struct{M_1} \iso[p] \Struct{M_2}$
  en volgt de stelling uit \olref[pis]{thm:p-isom1}.

  Eerst de heen-regel voor $\PIso{I}$. Neem $p \in
  \PIso{I}$ en schrijf $p(a_i) = b_i$ voor $i = 1$, \dots,~$n$.
  We mogen de elementen zo nummeren dat $a_1 <_1 a_2 <_1 \cdots <_1
  a_n$. Neem nu een $a \in \Domain{M_1}$. We zoeken op de volgende
  manier een passende $b \in \Domain{M_2}$:
  \begin{enumerate}
  \item als $a <_1 a_1$, kies $b \in \Domain{M_2}$ zo dat $b <_2
    b_1$;
  \item als $a_n <_1 a$, kies $b \in \Domain{M_2}$ zo dat $b_n <_2 b$;
 \item als $a_i <_1 a <_1 a_{i+1}$ voor een bepaalde $i$, kies dan $b \in
   \Domain{M_2}$ zo dat $b_i <_2 b <_2 b_{i+1}$.
  \end{enumerate}
  Zo'n $b$ is altijd te vinden. Dat volgt precies uit het feit dat
  $\Struct{M_2}$ dicht is en geen eerste of laatste element heeft. Neem nu
  $q = p \cup \{ \langle a, b \rangle \}$. Dan is $q \in \PIso{I}$ de
  uitbreiding van $p$ die we zochten. Daarmee is de heen-regel bewezen.
  De weer-regel werkt op dezelfde manier, met de twee ordeningen omgewisseld.
  Dus $\Struct{M_1} \iso[p] \Struct{M_2}$; volgens
  \olref[pis]{thm:p-isom1} volgt $\Struct{M_1} \iso \Struct{M_2}$.
\end{proof}

\begin{prob}
  Maak het bewijs van \olref[mod][bas][dlo]{thm:cantorQ} af. Controleer
  daarvoor dat $\PIso{I}$ ook aan de weer-regel voldoet.
\end{prob}

\begin{rem}
  Neem om het even welke !!{enumerable} dichte lineaire ordening zonder
  eindpunten, en noem die $\Struct{S}$. Volgens \olref{thm:cantorQ} geldt
  $\Struct{S} \iso \Struct{Q}$. Hier is
  $\Struct{Q} = (\Rat, <)$ de !!{enumerable} dichte lineaire ordening met
  de verzameling $\Rat$ van de rationale getallen als domein. Kijk nu opnieuw
  naar de !!{structure}~$\Struct{R} = (\Real, <)$ uit
  \olref[thm]{remark:R}. We zagen dat er !!a{enumerable}
  !!{structure}~$\Struct{S}$ bestaat waarvoor $\Struct{R} \elemequiv
  \Struct{S}$. Die $\Struct{S}$ is !!a{enumerable} dichte lineaire ordening
  zonder eindpunten. Daarom is zij isomorf (en dus ook elementair equivalent)
  met de !!{structure}~$\Struct{Q}$. Elementaire equivalentie
  is transitief. Daarom $\Struct{R} \elemequiv \Struct{Q}$. (We hadden dit
  ook rechtstreeks kunnen bewijzen: met hetzelfde heen-en-weerargument krijg
  je $\Struct{R} \iso[p] \Struct{Q}$.)
\end{rem}
\end{document}
